\chapter{Baker-Campbell-Hausdorff and Magnus Expansions}

\section{Introduction: When Exponentials Don't Commute}

\lettrine{T}{he} exponential function is the mathematician's Swiss Army knife—it solves differential equations, diagonalizes matrices, and connects algebra to geometry. But when the arguments don't commute, the simple identity $e^{A}e^{B} = e^{A+B}$ shatters into a beautiful cascade of commutators.

\begin{curiosity}{The Non-Commutative Landscape}{noncommutative}
In quantum mechanics, operators $A$ and $B$ rarely commute. The difference between $e^{A}e^{B}$ and $e^{B}e^{A}$ encodes physical information about uncertainty relations and measurement disturbances.
\end{curiosity}

\section{The Baker-Campbell-Hausdorff Formula}
% Content to be added in Task 2

\section{The Magnus Expansion}
% Content to be added in Task 3

\section{Connections and Applications}
% Content to be added in Task 4

\section{Computational Aspects}
% Content to be added in Task 5